\subsection{A complete intrinsic family with full geometric pro-\texorpdfstring{$3$}{3} lifts}

\begin{statement}{E.34}{Proved}{Every cubic abelian class unramified outside three lifts fully}
Let $N$ be the free pro-$3$ geometric fundamental group of the
plus Fermat cubic with its nine cusps removed, of rank ten.
Use the rational tangential section above the base tangent
$t=0$, $dt=1$, specified by
$s=-x/z$, $t=s^3$, $y/z=-(1-s^3)^{1/3}$.
Write $V=N^{\mathrm{ab}}=T_3J_{E,D}$ and
$G_S=G_{\mathbb Q,\{3,\infty\}}$.
The genuine action on $N$ factors through $G_S$.
Here a section is called unramified outside three when its
geometric cocycle relative to this tangent factors through $G_S$.

The entire group of such geometrically abelian section classes is
\[
 H^1_{\mathrm{cts}}(G_S,V)
       \simeq \mathbb Z_3^6\oplus(\mathbb Z/3)^3 .
 \tag{E34.1}
\]
Every one of these classes lifts to a section of
$N\rtimes G_{\mathbb Q}$ whose geometric cocycle also factors
through $G_S$. In particular all of them lift to the intrinsic
metabelian, intrinsic class-two and marked-depth quotients.
The actual coordinate-unit map sends (E34.1) onto
\[
 \{(a\kappa_3,b\kappa_3):a,b\in\mathbb Z_3\}
       \simeq\mathbb Z_3^2
 \tag{E34.2}
\]
in the toric minimum. Its kernel is
$\mathbb Z_3^4\oplus(\mathbb Z/3)^3$, and every abelian class
in each of these fibers has a full geometric pro-$3$ lift.
These are assertions about the specified unramified family,
not about all the ramified classes of E.24.

\emph{Manuscript: \S51.16, folios 733v--736v, scans 172--178,
(371)--(379); this is an intrinsic rank-ten geometric pro-$3$
family compared through the actual coordinate-unit map.}
\end{statement}
\begin{proof}
\emph{The actual action and its relative pro-$3$ quotient.}
The projective Fermat cubic over $\mathbb Z[1/3]$ is smooth:
its three partial derivatives cannot vanish simultaneously.
Each of its three coordinate-cusp schemes is finite etale of
degree three over that ring, and they are pairwise disjoint.
The specified cusp and parameter are integral away from three.
For the based comparison, the displayed Puiseux formulas send
$t^{1/m}$ to $s^{3/m}$ and commute with coefficientwise
Galois action. The constant cubic roots are $-1$ and have
compatible rational $3$-power roots. Thus $N$ is the actual
index-nine subgroup of the base free pro-$3$ group $F$ and
its action is the restriction of the standard based action
on $F$, with no additional pro-$3$ Kummer twist.

The based unramifiedness used here is stated explicitly in
Hain--Matsumoto, \S8.5, printed p.~28 of the opened version
\cite{E_HainMatsumoto}: take $X_N=\mathbb P^1-\{0,\mu_N,\infty\}$,
$N=1$ and $\ell=3$ there. Smooth base change gives the
\emph{genuine} action on the fundamental group with tangent
$\overrightarrow{01}$ over $G_S$. Restricting to the above
subgroup proves the asserted unramifiedness for $N$.

Put $K=\mathbb Q(\zeta)$ and let $P$ be the maximal pro-$3$
quotient of $G_{K,\{3\}}$. Its defining kernel is stable under
$G_S$, so the relative quotient
\[
 \Gamma=G_S/\ker(G_{K,\{3\}}\longrightarrow P)
 \quad\text{fits into}\quad
 1\longrightarrow P\longrightarrow\Gamma
       \longrightarrow C_2\longrightarrow1 .
 \tag{E34.3}
\]
Complex conjugation splits this sequence.
Over $K$ the action on $F/\Phi(F)\simeq\mathbb F_3^2$ is
trivial, because $\mu_3\subset K$.
The kernel of
$\operatorname{Aut}(F)\to\operatorname{GL}(F/\Phi(F))$
is pro-$3$. One way to see this is to use the characteristic
lower exponent-$3$ series: successive automorphism kernels
embed in homomorphism groups between its finite
$\mathbb F_3$-graded pieces and are $3$-groups; the finite
Burnside basis theorem starts the induction.
The inverse limit is the asserted pro-$3$ kernel.
The action of $G_{K,\{3\}}$ on $F$, and hence on $N$, therefore
factors through $P$. The full action factors through (E34.3).

The same factorization holds for every $V$-cocycle on $G_S$.
Indeed the map $g\mapsto(\alpha(g),\rho(g))$ takes
$G_{K,\{3\}}$ into the pro-$3$ group
$V\rtimes\rho(G_{K,\{3\}})$. By the defining universal property
of $P$ it kills the kernel in (E34.3).
This proves that both the action and all the cocycles under
consideration factor through the same actual quotient.

\emph{The free Galois kernel and vanishing of obstructions.}
The norm-Euclidean ring $\mathbb Z[\zeta]$ is a PID, its units
are $\mu_6$, and it has the single prime
$(\lambda)$, $\lambda=1-\zeta$, above three.
Kummer theory gives
\[
 H^1(G_{K,\{3\}},\mathbb F_3)
  =\langle\zeta,\lambda\rangle\simeq\mathbb F_3^2 .
 \tag{E34.4}
\]
For completeness, an element defining a cubic extension
unramified away from three has valuation divisible by three
at every other prime. Unique factorization then writes its
class as $\zeta^a\lambda^b$ modulo cubes.
Conversely these classes define extensions unramified away
from three. Their independence follows first from the
$\lambda$ valuation and then from the unit group.

The global Euler characteristic of
Milne, I.5.1 \cite{E_MilneADT}, for this imaginary quadratic
field and the trivial module $\mathbb F_3$ is $1/3$.
The groups in degrees zero and one have orders three and
nine, respectively, by (E34.4). Therefore
\[
 H^2(G_{K,\{3\}},\mathbb F_3)=0 .
 \tag{E34.5}
\]
Inflation embeds $H^2(P,\mathbb F_3)$ into this group.
To verify that point, let $J$ be the kernel defining $P$.
An invariant nonzero character $J\to\mathbb F_3$ would
produce a quotient of $G_{K,\{3\}}$ which is an extension
of $P$ by a central group of order three. That quotient is
pro-$3$, contradicting maximality of $P$.
The five-term sequence therefore proves the claimed injection.

These calculations prove that $P$ is free pro-$3$ of rank two.
Here is a direct explanation of the criterion being used.
The dimension in (E34.4) gives a minimal surjection from a
free pro-$3$ group on two generators to $P$.
The vanishing of $H^2(P,\mathbb F_3)$ gives vanishing in
degree two for every finite $3$-primary $P$-module:
each such module has a filtration with trivial
$\mathbb F_3$ factors, since its action image is a finite
$3$-group. Every finite embedding problem with $3$-group
kernel is then soluble. Indeed use a nonzero characteristic
elementary abelian subgroup of the center of its kernel,
solve the quotient problem by induction, and split the
remaining abelian extension using this $H^2$ vanishing.
For a characteristic open subgroup $U$ of the minimal free
group $F_0$ and its surjection $f:F_0\to P$, the precise
finite-kernel problem is the pullback of
$F_0/U\to P/f(U)$ along $P\to P/f(U)$.
Its solution gives the required homomorphism $P\to F_0/U$.
Apply this to all those finite quotients.
There are finitely many homomorphisms to each finite target,
because $P$ has two generators; compactness gives a compatible
section of the minimal surjection. Its resulting idempotent
endomorphism of the free group is the identity on the
Frattini quotient, hence is an automorphism by the pro-$3$
basis theorem. An idempotent automorphism is the identity,
so the minimal surjection is an isomorphism.

In particular $P$ has $3$-cohomological dimension one.
The Hochschild--Serre sequence for (E34.3), with averaging
over $C_2$ (two is invertible on every $3$-primary module),
gives
\[
 H^2(\Gamma,M)=0
 \quad\text{for every finite $3$-primary $\Gamma$-module }M.
 \tag{E34.6}
\]
The same conclusion holds over $G_S$ for every finite
module whose restriction to $G_{K,\{3\}}$ has $3$-group
action image: filter the restriction by trivial
$\mathbb F_3$ factors, apply (E34.5), and again take
$C_2$ invariants.

Now take any cocycle $\alpha:\Gamma\to V$.
Choose a decreasing characteristic series of $N$ with
finite $3$-group quotients $N_j$ and intersection one.
At each level, pull the surjection
$N_j\rtimes\Gamma\to N_j^{\mathrm{ab}}\rtimes\Gamma$
back along the reduction of $\alpha$.
Its kernel is a finite $3$-group. The same central-subgroup
induction using (E34.6) supplies a section.
There are finitely many choices at every level:
$\Gamma$ has the two generators of $P$ and complex
conjugation, and the geometric target is finite.
The compatible-section sets form a nonempty inverse limit.
Explicitly, every finite collection of levels is supplied
by a section at their largest common refinement; compactness
of the product of the finite sets gives a coherent choice.
The limit is a continuous $N$-cocycle lifting $\alpha$.
Inflating along $G_{\mathbb Q}\to G_S\to\Gamma$ gives
the asserted section of $N\rtimes G_{\mathbb Q}$.
This constructs lifts in the full geometric pro-$3$ group,
not merely compatible names for degree-two obstruction classes.

\emph{Counting the entire unramified abelian family.}
The genuine action on $V/3^nV$ has $3$-group image over
$G_{K,\{3\}}$, by the action factorization already proved.
The last assertion following (E34.6) gives
$H^2(G_S,V/3^nV)=0$ for every $n$.
These finite-coefficient degree-one groups are finite
(Milne I.4.15); hence the inverse-limit sequence also gives
$H^2_{\mathrm{cts}}(G_S,V)=0$.
Consequently
\[
 H^1_{\mathrm{cts}}(G_S,V)/3^n
       \simeq H^1(G_S,V/3^nV).
 \tag{E34.7}
\]

Complex conjugation has four plus eigenvalues and six minus
eigenvalues on $V$. To compute them, it fixes the three
rational cusps and interchanges three pairs among the other
six cusps. Its permutation action on degree-zero cusp
divisors therefore has multiplicities five and three.
The toric Tate twist reverses their signs, giving
multiplicities three and five on the rank-eight toric part.
On $T_3E$, its order is two and its determinant is $-1$ by
the Weil pairing, so its two eigenvalues are $1,-1$.
Adding the exact toric and elliptic pieces gives the claimed
multiplicities. Since two is a unit in $\mathbb Z_3$, the
same ranks hold for the integral eigensublattices.

By E.24 the rational $3$-primary torsion of $G=J_{E,D}$ is
$(\mathbb Z/3)^3$. Thus
$H^0(G_S,V/3^nV)=G(\mathbb Q)[3^n]$ has order $27$
for every $n\ge1$.
Milne's Euler formula I.5.1, now over $\mathbb Q$, has
real factor $3^{4n}/3^{10n}=3^{-6n}$.
Together with the preceding $H^2$ vanishing it yields
\[
 |H^1(G_S,V/3^nV)|=3^{6n+3}.
 \tag{E34.8}
\]
The continuous degree-one group is a finitely generated
$\mathbb Z_3$-module: its cocycles factor through the
three-generated group $\Gamma$ and form a closed
submodule of $V^3$, modulo the closed coboundary image.
Equations (E34.7)--(E34.8) give rank six and torsion of
order $27$ and exponent three. Equivalently that torsion is
$H^0(G_{\mathbb Q},V\otimes\mathbb Q_3/V)
=G(\mathbb Q)[3^\infty]$.
This proves (E34.1).

Finally the algebraic retraction of E.24 gives a split map
of these Galois modules to $\mathbb Z_3(1)^2$.
Kummer theory over $\mathbb Z[1/3]$ gives
$H^1(G_S,\mathbb Z_3(1))=\mathbb Z_3\kappa_3$.
Its rank-two target is torsion-free, and the splitting
therefore proves (E34.2) and the stated rank-four kernel
with its three torsion factors.
The change from the origin at $O$ in E.24 to the present
rational cusp is also unramified outside three. The framed
cusp-difference relation (E24.6) involves only the units
$\lambda$, powers of $3$ and signs, so this framed point
extends over $\mathbb Z[1/3]$ and its Kummer cocycle is
unramified there. Equivalently both integral tangential
fiber functors give the same unramified comparison.
Both origins have zero toric pro-$3$ class.
Thus changing origins preserves the unramified family
and the displayed linear toric map.
The matrix and order calculations are recorded in
\texttt{e\_intrinsic\_positive.py}; the Galois vanishing
and lifting arguments above establish the infinite assertions.
\end{proof}
